\documentclass[UTF8,fontset=fandol,11pt,a4paper]{ctexart}
\usepackage[a4paper,margin=1.75cm,headheight=15pt]{geometry}
\usepackage{amsmath,amssymb,mathtools}
\usepackage{booktabs,tabularx,array,longtable}
\usepackage{enumitem}
\usepackage[most]{tcolorbox}
\usepackage{tikz}
\usetikzlibrary{arrows.meta,positioning,matrix,calc}
\usepackage{xcolor,fancyhdr,hyperref,microtype,lastpage}

\newcolumntype{Y}{>{\raggedright\arraybackslash}X}
\newcolumntype{L}[1]{>{\raggedright\arraybackslash}p{#1}}
\newcolumntype{C}[1]{>{\centering\arraybackslash}p{#1}}
\setlength{\parindent}{2em}
\setlength{\parskip}{0.34em}
\setlength{\emergencystretch}{3em}
\renewcommand{\arraystretch}{1.2}
\setlength{\tabcolsep}{3pt}
\setlist[itemize]{itemsep=0.18em,topsep=0.35em,leftmargin=2em}
\setlist[enumerate]{itemsep=0.18em,topsep=0.35em,leftmargin=2em}

\definecolor{deep}{RGB}{38,58,72}
\definecolor{soft}{RGB}{244,247,248}
\definecolor{greenish}{RGB}{52,91,74}
\definecolor{warnc}{RGB}{136,61,58}
\definecolor{purple}{RGB}{84,72,112}
\definecolor{rulegray}{RGB}{110,116,120}

\hypersetup{colorlinks=true,linkcolor=deep,urlcolor=purple,pdftitle={CO-04 流水线与指令级并行}}
\pagestyle{fancy}
\fancyhf{}
\lhead{\small\color{deep}\textbf{408 · 计算机组成原理 CO-04}}
\rhead{\small\color{deep}Need · Ready · Overlap · Commit}
\cfoot{\small 第 \thepage\ 页 / 共 \pageref{LastPage} 页}
\renewcommand{\headrulewidth}{0.4pt}

\newtcolorbox{corebox}[1]{breakable,colback=soft,colframe=deep,title=\textbf{#1},boxrule=0.8pt,arc=2mm}
\newtcolorbox{methodbox}[1]{breakable,colback=white,colframe=greenish,title=\textbf{#1},boxrule=0.7pt,arc=1.5mm}
\newtcolorbox{warnbox}[1]{breakable,colback=white,colframe=warnc,title=\textbf{#1},boxrule=0.7pt,arc=1.5mm}
\newtcolorbox{examplebox}[1]{breakable,colback=white,colframe=purple,title=\textbf{#1},boxrule=0.7pt,arc=1.5mm}
\newcommand{\kw}[1]{\textbf{\color{deep}#1}}
\newcommand{\code}[1]{\texttt{#1}}

\tikzset{
  co/.style={draw=deep,rounded corners=1.5mm,text width=2.5cm,minimum height=0.9cm,align=center,fill=soft,font=\small,inner sep=3pt},
  state/.style={co,double,double distance=1.2pt,fill=white},
  flow/.style={-{Latex[length=2mm]},thick,draw=deep},
  ref/.style={-{Latex[length=2mm]},dashed,draw=rulegray}
}

\title{\vspace{-1.1cm}\textbf{流水线与指令级并行}\\[0.4em]
\large 多条指令怎样合法重叠而不破坏顺序语义}
\author{408 · 计算机组成原理 Topic CO-04}
\date{}

\begin{document}
\maketitle

\begin{corebox}{母问题：怎样重叠执行，同时保持与 ISA 顺序轨迹等价？}
CO-03 已给出一条指令内部的值依赖和提交边界。CO-04 把多条指令放到同一时间轴上，生成链是
\[
\boxed{
\text{Instruction Dependencies}
\to\text{Stage / Resource Occupancy}
\to\text{Need--Ready Constraints}
\to\text{Legal Overlap}
\to\text{Forward / Stall / Flush}
\to\text{Precise Commit}}
\]
箭头表示约束求解顺序：先承认程序依赖和资源事实，再判断哪些重叠合法；不是先背“停几拍”，再倒推理由。
\end{corebox}

\begin{methodbox}{本册定位与 Owner 边界}
\begin{itemize}
  \item \kw{Owns：}stage timing、latency/throughput、流水寄存器、结构/数据/控制 hazard、need/ready、forwarding、stall、flush、CPI、精确异常与 408 范围内 ILP 扩展。
  \item \kw{Uses：}CO-02 的 ISA 顺序语义；CO-03 的单指令值依赖、资源和提交；CO-06/07 的 miss/fault 延迟接口。
  \item \kw{Stops：}不重讲单指令微操作；不拥有 Cache/TLB 状态机；多核一致性、完整乱序执行和编译器后端只作 Extension。
\end{itemize}
\end{methodbox}

\tableofcontents
\newpage

\section{流水化从哪里生成}

\subsection{朴素方案与失败}
若一条指令依次经历多个功能阶段，非流水实现会让当前未使用的部件闲置。把长组合路径切成 $k$ 段，并在段间加入流水寄存器后，不同指令可以占用不同阶段。

设各段组合延迟为 $T_i$，流水寄存器开销为 $T_{reg}$，则
\[
T_{clk}\ge \max_i T_i+T_{reg}.
\]
理想情况下 $n$ 条指令总时间为
\[
T_{pipe}=(k+n-1)T_{clk},
\]
但这只是无 hazard、无 miss、段延迟均衡的成本模型。

\subsection{吞吐率与延迟}
流水线主要提高稳态吞吐率，不保证降低单条指令 latency。更深的流水可能缩短 $T_{clk}$，却增加流水寄存器开销、填充/排空、分支惩罚和旁路复杂度。因此比较性能必须回到
\[
T_{CPU}=IC\times CPI\times T_{clk}.
\]

\section{时间模型：Produced、Ready、Need 与 Commit}

对一个值 $v$，必须至少标四个时刻：
\begin{itemize}
  \item \kw{Produced：}功能部件已经计算出 $v$；
  \item \kw{Ready：}从某条实际路径看，消费者能够获得 $v$；
  \item \kw{Need：}消费者所在阶段必须使用 $v$；
  \item \kw{Commit：}$v$ 或其他效果写入软件可见架构状态。
\end{itemize}

RAW 依赖是否需要停顿取决于
\[
t_{ready}(producer,path)\le t_{need}(consumer).
\]
若不等式失败，必须改变时间（stall）或改变路径（增加 forwarding）；若值尚未 produced，任何旁路都无能为力。

\begin{warnbox}{Ready 不等于 Commit}
forwarding 可以让 ALU 结果在写回前被下一条指令使用，但寄存器堆此时可能仍是旧值。可用时刻回答“谁能拿到”，提交时刻回答“软件可见状态何时改变”。
\end{warnbox}

\section{Stage 与时空图：先声明模型}

经典五级教学模型常写为 IF、ID、EX、MEM、WB，但停顿数依赖以下假设：
\begin{itemize}
  \item 各 stage 的真实职责与输入/输出；
  \item 寄存器堆同周期写/读的先后关系；
  \item 指令存储与数据存储是否共享端口；
  \item 分支在哪一级比较并产生目标；
  \item forwarding 的起点和终点；
  \item memory 是否固定一拍返回。
\end{itemize}

\begin{center}
\small
\begin{tabularx}{\linewidth}{C{1.5cm}YYYYY}
\toprule
周期 & IF & ID & EX & MEM & WB\\
\midrule
1 & $I_1$ & & & &\\
2 & $I_2$ & $I_1$ & & &\\
3 & $I_3$ & $I_2$ & $I_1$ & &\\
4 & $I_4$ & $I_3$ & $I_2$ & $I_1$ &\\
5 & $I_5$ & $I_4$ & $I_3$ & $I_2$ & $I_1$\\
\bottomrule
\end{tabularx}
\end{center}

时空图不是装饰：每个格子同时表示 stage、资源占用和流水寄存器边界。发生 stall 时必须说明哪些前级保持、哪个控制包被清零为 bubble、哪些后级继续推进。

\section{三类 Hazard：三种约束冲突}

\subsection{结构冲突}
同一周期两项工作需要同一不可并用资源，例如 IF 与 MEM 共用单端口存储器。解决路径只有三类：复制/增加端口，以空间换吞吐；错开调度，以周期换资源；或改变组织。是否冲突由题设资源图决定，不能由“无数据依赖”推出无冲突。

\subsection{数据冲突}
\begin{itemize}
  \item \kw{RAW：}消费者读生产者将写的新值，是真依赖；
  \item \kw{WAR：}较年轻指令写得过早，破坏较老指令尚未完成的读；
  \item \kw{WAW：}两个写同一目标的提交次序颠倒。
\end{itemize}

在单发射、按序发射、按序读写的简单五级流水中，通常只有 RAW 形成实际 hazard；WAR/WAW 需要乱序读写、多发射或不同完成延迟才可能出现。这个结论必须绑定模型，不能写成所有流水线的定律。

\subsection{控制冲突}
分支方向/目标尚未确定时，前级已经取入后续指令。等待、静态/动态预测、延迟槽和尽早决策都在交换平均停顿、硬件/编译器复杂度和错误路径成本。无论策略如何，错误路径效果不能提交。

\section{三种控制动作：Forward、Stall、Flush}

\begin{center}
\small
\begin{tabularx}{\linewidth}{L{2.0cm}Y Y Y}
\toprule
动作 & 保留/推进什么 & 阻止什么 & 使用条件\\
\midrule
Forward & 正常推进并改用旁路值 & 消费旧值 & 值已经 produced，且实际路径可达 need 端\\
Stall & 保持 PC/前级状态或局部队列 & 未准备好的消费者前进 & need 早于所有 ready 路径，或资源冲突\\
Flush & 保留正确较老状态 & 错路/异常后的年轻效果提交 & 分支误判、异常、重定向\\
\bottomrule
\end{tabularx}
\end{center}

forwarding 不是“打开一个总旁路”。ALU $\to$ EX、load data $\to$ EX、result $\to$ branch compare、store data $\to$ MEM 是不同路径；题目没给路径，就不能假设存在。

\section{贯穿母例：LOAD 后立即 ADD}

\[
I_1:\ \code{LOAD R1,0(R2)},\qquad
I_2:\ \code{ADD R3,R1,R4}.
\]

在经典五级模型中：$I_1$ 的 EA 在 EX 产生，数据通常在 MEM 末尾产生；$I_2$ 在自己的 EX 开始需要 $R1$。若相邻发射，数据的 ready 晚于 need，必须至少让消费者后移，直到存在可达旁路。

\begin{center}
\small
\begin{tabularx}{\linewidth}{C{1.1cm}YYYYYYY}
\toprule
 & 1 & 2 & 3 & 4 & 5 & 6 & 7\\
\midrule
$I_1$ & IF & ID & EX & MEM & WB & &\\
$I_2$ & & IF & ID & stall & EX & MEM & WB\\
\bottomrule
\end{tabularx}
\end{center}

该“一拍”只在上述 stage、memory latency、寄存器时序和 MEM-to-EX 旁路假设下成立。若分级更深、memory 多周期或旁路端点不同，必须重新比较 need/ready。

\begin{warnbox}{母例故意暴露的错误路径}
“有 forwarding，所以永不停顿”忽略了值尚未产生；“LOAD 在 EX 算出了结果”把 EA 当成 load data；“I1 已经转发，所以已经完成”把 ready 当 commit。
\end{warnbox}

\section{控制冒险与预测状态机}

\subsection{惩罚从决策位置生成}
若分支在第 $d$ 级末尾才确定，错误路径已经进入的年轻指令数量由前端组织和预测时机决定。penalty 不是“分支固定几拍”，必须从决策级、已推进阶段和 flush 范围数出。

平均 CPI 贡献可写为
\[
CPI_{branch}=f_{branch}\,r_{mispredict}\,P_{mispredict},
\]
其中 $P_{mispredict}$ 必须是额外周期；若题目给的是总分支周期，需先剥离理想基线，避免重复计数。

\subsection{两位饱和计数器}
两位预测器是四状态 FSM，单次相反结果只使状态向另一方向移动一步。题目必须给初态、实际 T/NT 序列和采用高位/状态预测规则；循环“通常只错一次”不是无条件定律，初态和连续多轮会改变结果。

\section{CPI 与总时间：从时空图计数}

实际 CPI 可以分解为
\[
CPI_{actual}=CPI_{ideal}
+CPI_{struct}+CPI_{data}+CPI_{control}+CPI_{memory}+\cdots.
\]
分项只有在互斥或题设明确时才能直接相加；若一个 memory miss 同时阻止了后续依赖，不能把同一停顿重复归入两项。

有限指令序列的总周期应从时空图第一拍数到最后一条的规定完成/提交点。填充与排空已经包含在图中，不应再机械追加 $k-1$。

\section{精确异常与提交顺序}

精确状态要求：异常指令之前的较老指令效果已经成立；异常指令和更年轻指令的禁止副作用没有提交。顺序流水可以冻结/清零写使能并 flush；更复杂实现把 produced/finished 与 ordered commit 分开。

\[
\boxed{\text{execution may overlap}\quad\land\quad
\text{architectural trace remains recoverable}}
\]

这条不变量连接 CO-03 的单指令 Write Set、CO-04 的多指令年龄顺序和 X-B01 的异常入口。OS handler 不在本册展开。

\section{代价、边界与 ILP 扩展}

\begin{center}
\small
\begin{tabularx}{\linewidth}{L{2.2cm}Y Y Y}
\toprule
机制 & 换来的能力 & 新成本 & 最小失败条件\\
\midrule
更深流水 & 更短候选时钟、更多重叠 & 寄存器开销、分支/依赖周期数增加 & 最慢段/寄存器开销不降\\
Forwarding & 减少已产生值的 RAW stall & 比较器、MUX、布线和时序压力 & 值尚未产生或端点不可达\\
分支预测 & 降低平均控制停顿 & 预测状态、能耗、错误 flush & 低可预测性或高 penalty\\
多发射 & 峰值 IPC 大于 1 & 端口、依赖检查、带宽 & 指令不独立或资源不足\\
VLIW & 编译期显式并行、简化部分硬件 & 代码密度、调度和兼容性压力 & 动态延迟/并行不足\\
\bottomrule
\end{tabularx}
\end{center}

超标量、乱序、寄存器重命名、ROB、VLIW 和多核属于“同一依赖/资源/提交框架怎样扩张”的 Extension。408 主干只保留能解释 WAR/WAW、动态调度和 ordered commit 的最小接口，不把现代实现参数写成固定规律。

\section{Flynn 与多处理器：并行对象先分清}
流水线和超标量主要是在一条指令流内重叠；多处理器则增加执行上下文或指令流。按指令流/数据流分类：
\begin{center}\small
\begin{tabularx}{\linewidth}{L{2.0cm}L{2.2cm}YY}
\toprule
类别 & 指令流/数据流 & 典型结构 & 题面判据\\\midrule
SISD & 单/单 & 单核顺序处理器 & 一条指令处理一个数据序列\\
SIMD & 单/多 & 向量机、向量指令 & 同一控制对多个数据元素执行\\
MISD & 多/单 & 少见的专用容错结构 & 多条指令作用于同一数据流\\
MIMD & 多/多 & 多核、SMP、通用多处理器 & 不同核心可执行不同线程/程序\\
\bottomrule
\end{tabularx}\end{center}
向量处理属于 SIMD 的数据级并行；多核/SMP 属于 MIMD 的共享或分布式执行资源，不能把“多发射每拍多条”直接说成“多核”。硬件多线程则是在一个核心内保存多个线程上下文，在某线程等待或资源空闲时切换发射；它提高资源利用率，不自动增加物理 ALU 数量。

多处理器题仍沿本册三条约束检查：每个执行流的依赖与 ready/need、共享 Cache/总线/内存的资源冲突、最终架构提交是否符合各自 ISA 顺序。题目若进一步问 Cache 一致性、锁或线程调度，应交给相应跨册/OS Owner；本节只确定并行层次与硬件接口。

\begin{methodbox}{并行层级的停止条件}
先回答“增加的是数据元素、同一指令流中的发射宽度、线程上下文，还是物理核心/指令流”；再检查共享资源与提交边界。若无法指出新增并行对象，就不能仅凭“同时”判定属于 SIMD、超标量或多核。
\end{methodbox}

\section{五列概念边界}

\begin{center}
\small
\begin{tabularx}{\linewidth}{L{1.8cm}C{0.35cm}L{1.8cm}Y Y}
\toprule
概念 A & $\neq$ & 概念 B & 真正区别与题目信号 & 混淆后果\\
\midrule
Latency & $\neq$ & Throughput & 单条穿越时间 vs 单位时间完成数 & 误称流水必降单条延迟\\
依赖 & $\neq$ & Hazard & 程序语义关系 vs 当前时间/资源下的冲突 & 把所有 RAW 都算固定停顿\\
Produced & $\neq$ & Ready & 已算出 vs 通过实际路径可被消费者拿到 & 假设不存在的旁路\\
Ready & $\neq$ & Commit & 消费者可用 vs 架构状态已更新 & 精确异常和寄存器值判断错误\\
Stall & $\neq$ & Flush & 等待后继续同一路径 vs 作废错误/年轻状态 & 保留了不应提交的控制包\\
ILP & $\neq$ & 多核并行 & 单指令流内部重叠 vs 多核心/任务并行 & 把 OS 调度和 Cache 一致性塞入流水线题\\
\bottomrule
\end{tabularx}
\end{center}

\section{做题控制协议}

\begin{enumerate}
  \item 写 stage 定义、每级资源、寄存器读写时序、memory latency 和已有 forwarding；
  \item 逐条列源/目的寄存器和控制转移，标 producer、consumer；
  \item 对每个相关值标 produced、ready、need、commit；
  \item 先画无冲突时空图，再叠加结构、数据和控制约束；
  \item forwarding 必须写起点、终点和到达时刻；剩余 ready--need 差才用 stall；
  \item 分支写预测状态、决策级和 flush 范围；
  \item 从最终时空图统计总周期/CPI，不重复添加填充或同一停顿；
  \item 用年龄顺序和 Write Set 检查错误/异常路径没有提交。
\end{enumerate}

\begin{methodbox}{压缩成两问}
流水线题先问：\kw{这个值沿现有路径什么时候 ready，消费者什么时候 need？这个副作用在哪个年龄/时钟边界才允许 commit？}
\end{methodbox}

\section{全科坐标与跨册交接}
\begin{corebox}{Map Coordinate：Timing / Availability}
CO-04 接收 CO-03 的单指令路径，给每个值标出 produced、ready、consumer need 和 architectural commit；它向 CO-I01 输出合法重叠、forward/stall/flush 与精确提交顺序。阶段名称不能替代这些时刻。
\end{corebox}
吞吐量是稳定填满后的单位时间完成量，latency 是单条指令从进入到结果/提交的时间；两者和 clock rate、CPI、分支 penalty、miss stall 一起才形成 Cost。流水线加深或并行度增加不自动降低完整 CPU time。

\section{考纲映射与一页压缩}

\begin{center}
\small
\begin{tabularx}{\linewidth}{L{3.2cm}Y Y}
\toprule
传统考点 & 本册机制位置 & 下游接口\\
\midrule
流水线性能 & stage 延迟、填充/排空、CPU time & Rules 性能工具箱\\
结构/数据/控制冒险 & 资源、need/ready、Next PC & CO-03/CO-B01\\
转发、停顿、冲刷 & 三种控制动作与接口条件 & 具体时空图题\\
分支预测 & 预测 FSM、决策级、mispredict penalty & CO-I01\\
超标量/乱序/VLIW & ILP Extension 与按序提交边界 & 不扩张为多核 Topic\\
\bottomrule
\end{tabularx}
\end{center}

\begin{corebox}{一句话}
流水线把多条指令放到同一时间轴上：先用 need/ready 和资源约束判断哪些重叠合法，再用 forward、stall、flush 修复冲突，最终只允许符合年龄顺序和 ISA 的效果提交。
\end{corebox}

自测：为什么 forwarding 不能消灭所有 RAW？为什么 WAR/WAW 在简单五级顺序流水中通常不形成 hazard？为什么更深流水不保证更快？为什么分支 penalty 必须从决策级和 flush 范围数出？

\section{来源处理}

本册吸收归档《指令流水线》和原 Topic README 的 stage、hazard、forwarding、分支与性能主干。固定五级停顿数、特定寄存器半拍规则、现代 CPU 深度/宽度和“所有 RISC/CISC 性能相当”等实现结论均未升级为稳定定律；Flynn、多核、完整乱序和编译器后端只保留 Extension 指针。

\end{document}
